\documentclass[10pt]{article}

\usepackage[a4paper,margin=18mm,headheight=14pt]{geometry}
\usepackage[T1]{fontenc}
\usepackage{lmodern}
\usepackage{microtype}
\usepackage{amsmath,amssymb,mathtools,bm}
\usepackage{booktabs,tabularx,array}
\usepackage{enumitem}
\usepackage{xcolor}
\usepackage[
  colorlinks=true,
  linkcolor=blue!45!black,
  urlcolor=blue!45!black
]{hyperref}
\usepackage[capitalise,nameinlink,noabbrev]{cleveref}
\usepackage{fancyhdr}

\definecolor{maincolor}{HTML}{17365D}
\setlength{\parindent}{0pt}
\setlength{\parskip}{0.45\baselineskip}
\setlist{nosep,leftmargin=*}
\setcounter{secnumdepth}{2}
\setcounter{tocdepth}{1}
\renewcommand{\arraystretch}{1.12}

\pagestyle{fancy}
\fancyhf{}
\fancyhead[L]{\small TimeTensors forecasting benchmark}
\fancyhead[R]{\small Gaspard Berthelier}
\fancyfoot[C]{\thepage}
\renewcommand{\headrulewidth}{0.35pt}
\fancypagestyle{plain}{%
  \fancyhf{}
  \fancyhead[L]{\small TimeTensors forecasting benchmark}
  \fancyhead[R]{\small Gaspard Berthelier}
  \fancyfoot[C]{\thepage}
  \renewcommand{\headrulewidth}{0.35pt}%
}
\hypersetup{
  pdftitle={TimeTensors: Forecasting Benchmark and Experiments},
  pdfauthor={Gaspard Berthelier}
}

\newcommand{\reals}{\mathbb{R}}
\newcommand{\code}[1]{\texttt{\detokenize{#1}}}
\newcommand{\MSE}{\operatorname{MSE}}
\newcommand{\MAE}{\operatorname{MAE}}
\newcommand{\nMSE}{\operatorname{nMSE}}
\newcommand{\nMAE}{\operatorname{nMAE}}
\newcommand{\rMSE}{\operatorname{rMSE}}
\newcommand{\Wten}{\operatorname{W10}}

\title{\vspace{-1.5em}\textbf{TimeTensors: Forecasting Benchmark}\\[2mm]\large Experiment Guideline\vspace{-0.4em}}
\author{Gaspard Berthelier}
\date{7 August 2026}

\begin{document}
\maketitle
\thispagestyle{empty}

\begin{abstract}
TimeTensors studies how data selection, sampling, normalization, optimization
loss, model class, and training scope affect time-series forecasts. This
document defines the forecasting task, the seven implemented experiment
families, and the exact common protocol used by their launchers.
\end{abstract}

\tableofcontents

\section{Time-series forecasting task}
\label{sec:task}

Let $\mathcal U=\{0,\ldots,U-1\}$ index series (called \emph{users} in the
code), let $d\in\{1,\ldots,D\}$ index variables within a series, and let
$t\in\{0,\ldots,T-1\}$ index ordered dates. The complete panel is
\[
  Z=(z_{u,d,t})\in\reals^{U\times D\times T}.
\]
The benchmark CSV loader treats each non-date column as one series, hence all
publication datasets use $D=1$; the implementation nevertheless retains the
general $D$-variate tensor contract.

Fix a lookback length $L\geq1$ and a forecast horizon $H\geq1$. A query date
$t$ is the last observed date. Its input and target are
\begin{align}
  X_{t,u}
  &=Z_{u,:,\,(t-L,t]}
    =(z_{u,:,t-L+1},\ldots,z_{u,:,t})\in\reals^{D\times L},\\
  Y_{t,u}
  &=Z_{u,:,\,(t,t+H]}
    =(z_{u,:,t+1},\ldots,z_{u,:,t+H})\in\reals^{D\times H}.
\end{align}
Thus $L$ is the number of available past values and $H$ the number of future
values predicted at once. A forecasting model is a map
\[
  F_\theta:\reals^{D\times L}\longrightarrow\reals^{D\times H},
  \qquad \widehat Y_{t,u}=F_\theta(X_{t,u}),
\]
where $\theta$ contains trainable parameters for learned models and is empty or
frozen for persistence and Chronos-2.

The dates are split chronologically into target-date blocks
$S_{\mathrm{tr}},S_{\mathrm{va}},S_{\mathrm{te}}$ in proportions
$0.6/0.2/0.2$. For any block $S$, the admissible queries are
\[
  \mathcal T(S)=\{t:L-1\leq t\leq T-H-1,
  \ \{t+1,\ldots,t+H\}\subseteq S\}.
\]
The full target therefore belongs to its split, while the lookback may cross
the preceding split boundary. In particular, changing $L$ does not change the
validation or test target dates. The default training and evaluation pair sets
are $\mathcal A_S=\mathcal U\times\mathcal T(S)$, before the filters and strides
defined below.

For a lookback $X$, statistics are computed independently along its time axis:
\[
  \mu_X=\frac1L\sum_{\ell=1}^{L}X_{:,\ell},\qquad
  \sigma_X=\left(\frac1L\sum_{\ell=1}^{L}
  (X_{:,\ell}-\mu_X)^2\right)^{1/2},\qquad \varepsilon=10^{-8}.
\]
The population standard deviation is used. Model normalization is inverted
before the loss is evaluated, so predictions and targets passed to every loss
are in the original data space.

\section{Experiment 1: constant windows and users}
\label{sec:constants}

For the univariate benchmark, an accessible pair $(u,t)$ is valid when its
lookback is finite and non-constant:
\[
  a_{u,t}=\bm 1\!\left[
  X_{t,u}\text{ is finite and }\sigma_{X_{t,u}}>0
  \right].
\]
Window filtering removes only pairs for which $a_{u,t}=0$. User filtering is
more restrictive: for split $S$ it retains
\[
  \mathcal U_{S}^{\mathrm{safe}}
  =\{u\in\mathcal U:a_{u,t}=1\ \text{for every }t\in\mathcal T(S)\}.
\]
In the general $D$-variate implementation, window filtering accepts a
user--date pair if at least one variable is finite and non-constant, whereas
user filtering drops a user if any accessible variable-window is constant or
non-finite. These definitions coincide for the $D=1$ benchmark tensors.

The launcher compares exactly five policies in full and ultra: keep all pairs;
\path{remove_train_windows}, \path{remove_eval_windows}, or
\path{remove_all_windows}; and \path{drop_all_users}. Train and evaluation
interventions are deliberately separated: train-only filtering tests an
optimization effect, whereas evaluation filtering changes the measured
population. Curated source exclusions from each dataset's \path{config.json}
are applied before these dynamic policies and are not an experimental factor.

The family uses instance-normalized PatchTST in the full profile and adds
DLinear in ultra. Its MSE tables should be interpreted with the retained
numbers of users and windows, since lower error after evaluation filtering is
not by itself evidence of a better forecaster.

\section{Experiment 2: sampling and batch size}
\label{sec:sampling}

The sampler controls which accessible pairs form optimizer batches. With the
current default of one user per dataset item, its four modes are:
\begin{itemize}
  \item \code{random}: the single dataset item draws a query date and then one
    user uniformly at random;
  \item \code{dates}: item $i$ fixes the $i$th accessible date and draws one
    user at random;
  \item \code{individuals}: item $i$ fixes user $i$ and draws one accessible
    date at random;
  \item \code{all}: items deterministically enumerate every accessible
    $(u,t)$ pair in date-major order.
\end{itemize}
After filtering, let $N_m$ be the dataset length in mode $m$. Ignoring removed
pairs,
\[
  N_{\mathrm{random}}=1,\qquad
  N_{\mathrm{dates}}=|\mathcal T(S_{\mathrm{tr}})|,\qquad
  N_{\mathrm{individuals}}=U,\qquad
  N_{\mathrm{all}}=|\mathcal A_{S_{\mathrm{tr}}}|.
\]
For nominal batch size $B$ and configured epoch count $E$, the current training
loop performs
\[
  K_m=E\left\lceil\frac{N_m}{B}\right\rceil
\]
optimizer steps. Consequently, the full-profile value $E=10{,}000$ is exactly
10,000 optimizer steps only in random mode; the modes do not share a fixed
update budget. Moreover, random mode returns one actual window per batch with
the current dataset length, irrespective of nominal $B$. This behavior is part
of the present implementation and must be accounted for when interpreting the
batch-size sweep.

Full and ultra cross the four modes with $B\in\{64,256,1024\}$. Evaluation is
held to exhaustive pair enumeration with stride $H$; changing $B$ changes its
batching but not its query set or metric. The model uses instance normalization,
and the shared model axis is PatchTST in full and PatchTST plus DLinear in ultra.

\section{Experiment 3: normalization layers}
\label{sec:normalization}

The wrapper applies an invertible normalization $N$ before the base forecaster
$f_\theta$ and restores its output afterward:
\[
  \widehat Y=N^{-1}\!\left(f_\theta(N(X))\right).
\]
Let $(\mu_0,\sigma_0,a_0,b_0)$ be the scalar mean, standard deviation,
minimum, and maximum computed over all accessible training lookbacks. Let
$a_X=\min_\ell X_{:,\ell}$ and $b_X=\max_\ell X_{:,\ell}$. The sweep compares:
\begin{align*}
  N_{\mathrm{id}}(X)&=X,\\
  N_{\mathrm{std}}(X)&=\frac{X-\mu_0}{\sigma_0+\varepsilon},
  &N_{\mathrm{std}}^{-1}(Q)&=Q(\sigma_0+\varepsilon)+\mu_0,\\
  N_{\mathrm{mm}}(X)&=\frac{X-a_0}{b_0-a_0+\varepsilon},
  &N_{\mathrm{mm}}^{-1}(Q)&=Q(b_0-a_0+\varepsilon)+a_0,\\
  N_{\mathrm{imm},X}(X)&=\frac{X-a_X}{b_X-a_X+\varepsilon},
  &N_{\mathrm{imm},X}^{-1}(Q)&=Q(b_X-a_X+\varepsilon)+a_X,\\
  N_{\mathrm{inst},X}(X)&=\frac{X-\mu_X}{\sigma_X+\varepsilon},
  &N_{\mathrm{inst},X}^{-1}(Q)&=Q(\sigma_X+\varepsilon)+\mu_X.
\end{align*}
The code names are respectively \code{identity}, \code{standard},
\code{min-max}, \code{in-min-max}, and \code{instance}. The sixth method,
\code{revin}, adds one learnable scale and bias per variable:
\[
  N_{\mathrm{RevIN},X}(X)
  =\gamma\odot N_{\mathrm{inst},X}(X)+\beta,
  \quad
  N_{\mathrm{RevIN},X}^{-1}(Q)
  =\frac{Q-\beta}{\gamma+\varepsilon}
   (\sigma_X+\varepsilon)+\mu_X.
\]
Thus \code{instance} and \code{revin} use the same
\code{RevINNormalization} class, with affine parameters disabled and enabled,
respectively. All instance statistics are detached and cached from the current
lookback. The experiment fixes the training loss to nMSE and reports
inverse-transformed test MSE separately for each shared backbone.

\section{Experiment 4: reference orders of error}
\label{sec:reference}

This experiment establishes the error scale before architecture-specific
ablations. It compares three qualitatively different forecasters under the same
split and evaluation queries:
\begin{align*}
  \widehat Y^{\mathrm{pers}}_{t,u,:,h}&=X_{t,u,: ,L},
  &&h=1,\ldots,H,\\
  \widehat Y^{\mathrm{PatchTST}}_{t,u}&=F_{\widehat\theta}(X_{t,u}),\\
  \widehat Y^{\mathrm{Chronos}}_{t,u}&=F_{\theta_0}^{\mathrm{median}}(X_{t,u}).
\end{align*}
Persistence has no trainable parameters. PatchTST is trained from scratch with
instance normalization and nMSE. Chronos-2 is loaded from the local checkpoint,
kept frozen, uses identity normalization and no cross-learning, and returns the
median forecast quantile. Users with an accessible constant lookback are
dropped in both training and evaluation by default.

The family writes a combined test MSE table and a combined worst-user-decile
MSE table. The comparison is an order-of-error reference, not a claim that the
three methods have equal training data or compute.

\section{Experiment 5: training losses}
\label{sec:losses}

For one prediction, let $Q=DH$ and let $j$ index its $D\times H$ elements. The
five implemented training objectives are
\begin{align}
  \MSE(\widehat Y,Y)
  &=\frac1Q\sum_{j=1}^{Q}(\widehat Y_j-Y_j)^2,\\
  \MAE(\widehat Y,Y)
  &=\frac1Q\sum_{j=1}^{Q}|\widehat Y_j-Y_j|,\\
  \nMSE(\widehat Y,Y\mid X)
  &=\frac1Q\sum_{j=1}^{Q}
    \left(\frac{\widehat Y_j-Y_j}{\sigma_X+\varepsilon}\right)^2,\\
  \nMAE(\widehat Y,Y\mid X)
  &=\frac1Q\sum_{j=1}^{Q}
    \frac{|\widehat Y_j-Y_j|}{\sigma_X+\varepsilon},\\
  \rMSE(\widehat Y,Y\mid X)
  &=\frac1Q\sum_{j=1}^{Q}
    \left(\frac{\widehat Y_j-Y_j}{|\mu_X|+\varepsilon}\right)^2.
\end{align}
The scales broadcast from each sample and variable over its horizon. The code
name for the last objective is \code{relative_mse}. Because scaling is applied
to prediction and target before the base MSE/MAE, these equations match the
implementation exactly.

The sweep holds instance normalization fixed and changes only the training
criterion. Its tables compare every run using data-space test MSE, so training
curve magnitudes from different objectives are not directly comparable.
Complete run artifacts additionally retain MSE, nMSE, MAE, nMAE, relative MSE,
and per-user losses.

\section{Experiment 6: linear models}
\label{sec:linear}

Linear controls test how much performance follows from autoregressive structure
alone. The unrestricted trainable model is
\[
  \operatorname{vec}(\widehat Y)
  =W\operatorname{vec}(X)+b,
  \qquad W\in\reals^{DH\times DL},\quad b\in\reals^{DH}.
\]
The periodic model uses period $p=\min(L,168)$. With zero-based lag and horizon
indices, define
\[
  I_h=\{\ell\in\{0,\ldots,L-1\}:\ell\bmod p=(L+h)\bmod p\}.
\]
It fits one head per horizon,
\[
  \widehat Y_{:,h}=W_h\operatorname{vec}(X_{:,I_h})+b_h,
  \qquad h=0,\ldots,H-1,
\]
so each forecast step sees only past positions of the same phase. No limit is
placed on the number of previous cycles.

The \code{linear} and \code{periodic_linear} models are optimized by the common
PyTorch training loop. \code{sklinear} instead deterministically unrolls every
sampler-accessible training pair and fits scikit-learn multi-output ordinary
least squares with an intercept. Persistence is included as a parameter-free
reference. Each fitted linear method is evaluated with \code{identity},
\code{standard}, and non-affine \code{instance} normalization; the launcher
also saves its coefficient tensor and plots. The test profile narrows this to
\code{linear} and \code{sklinear} with instance normalization. All methods
share one MSE table.

\section{Experiment 7: central and per-user forecasting}
\label{sec:scope}

A central model fits one parameter vector using all training users. Per-user
training creates an independent model and output directory for each user in the
training split, restricts every available split to that user, and then merges
the resulting losses. For PatchTST this compares shared training with
specialization. For frozen Chronos-2, ultra mode instead compares central
inference with \code{cross_learning=true} against separate per-user inference
with \code{cross_learning=false}.

Let $n_u$ be the number of evaluated scalar errors for user $u$ and let
$\ell_{u,i}$ denote one such error. Define
\[
  \overline\ell_u=\frac1{n_u}\sum_{i=1}^{n_u}\ell_{u,i},\qquad
  M_{\mathrm{elem}}=\frac{\sum_u\sum_i\ell_{u,i}}{\sum_u n_u},\qquad
  M_{\mathrm{user}}=\frac1U\sum_u\overline\ell_u.
\]
For $q=\max(1,\lceil0.1U\rceil)$ and $\operatorname{Top}_q$ the users with
the $q$ largest means, the tail metric is
\[
  \Wten(\ell)=\frac1q\sum_{u\in\operatorname{Top}_q}\overline\ell_u.
\]
The current central artifacts store elementwise MSE, hence their reported
\code{mse} reduces to $M_{\mathrm{elem}}$. The per-user merger explicitly
overwrites \code{mse} with $M_{\mathrm{user}}$. The combined MSE table therefore
uses these two different weightings, while the W10 table uses the same
equal-user tail definition for both scopes. This is the implemented artifact
contract and should be stated whenever the combined table is interpreted.

\section{Experimental details}
\label{sec:details}

\subsection{Datasets and task grid}

Every CSV column other than \code{date} becomes one univariate series. The
configured publication datasets are listed in \cref{tab:datasets}; counts are
after the curated source exclusions in the current \code{config.json} files.

\begin{table}[htbp]
\centering
\small
\caption{Configured publication datasets. Solar and Weather are resampled to
hourly cadence during tensor construction.}
\label{tab:datasets}
\begin{tabularx}{\linewidth}{@{}lrrlX@{}}
\toprule
Dataset & Series & Raw cadence & Used cadence & Construction \\
\midrule
ETTh1 & 7 & 1 hour & 1 hour & All seven non-date variables. \\
Electricity & 312 & 1 hour & 1 hour & 321 raw series; nine configured exclusions. \\
Traffic & 862 & 1 hour & 1 hour & No configured source exclusion. \\
Solar & 137 & 10 minutes & 1 hour & Hourly sums. \\
Weather & 21 & 10 minutes & 1 hour & Hourly means. \\
Exchange Rate & 8 & 1 day & 1 day & No configured source exclusion. \\
\bottomrule
\end{tabularx}
\end{table}

Full and ultra use
\[
  (L,H)\in\{(168,24),(336,48),(504,168)\}
\]
and seeds $\{1,2,3\}$. Test mode uses only Electricity at $(504,168)$ with
seed 1. A dataset's sibling \code{config.json} is loaded during tensor
construction: portable top-level options are applied first, the
\code{timetensors} object overrides them, explicit run values override both,
and all \code{drop_users} lists are merged additively.

\subsection{Models}

\begin{table}[htbp]
\centering
\small
\caption{Models exercised by the seven publication launchers.}
\label{tab:models}
\begin{tabularx}{\linewidth}{@{}lXl@{}}
\toprule
Code name & Implemented forecast & Fit in this project \\
\midrule
\code{persistence} & Repeats the last lookback value for all $H$ steps. & None \\
\code{linear} & One affine map $\reals^{DL}\to\reals^{DH}$. & Adam \\
\code{periodic_linear} & One affine head per horizon using phase-matched lags, $p=\min(L,168)$. & Adam \\
\code{sklinear} & Multi-output scikit-learn \code{LinearRegression} over all accessible windows. & Closed-form OLS \\
\code{dlinear} & Moving-average decomposition (kernel 25) followed by separate per-variable seasonal and trend maps $\reals^L\to\reals^H$. & Adam \\
\code{patchtst} & Patches of length 16 and stride 8; three encoder layers, 16 heads, $d_{\rm model}=128$, $d_{\rm ff}=256$, learned positional encoding, and a flattened forecast head. & Adam \\
\code{chronos} & Local frozen Chronos-2 pipeline; the wrapper selects the median quantile and optionally enables cross-learning. & Frozen \\
\bottomrule
\end{tabularx}
\end{table}

PatchTST uses attention dropout 0, encoder dropout 0.2, feed-forward/head
dropout 0.2/0, GELU activations, and residual attention. DLinear and PatchTST
form the shared model axis: full uses PatchTST and ultra adds DLinear. The
reference and linear experiments use their fixed method lists instead. The
central/per-user family uses PatchTST in test and full, and adds Chronos-2 in
ultra.

\subsection{Optimization, evaluation, and outputs}
\thispagestyle{fancy}

Unless a family overrides a value, training uses Adam with PyTorch defaults,
learning rate $10^{-5}$, batch size 256, random training sampling with stride
1, exhaustive evaluation sampling with stride $H$, and no scheduler, early
stopping, or gradient clipping. Test runs $E=20$ epochs and validates/logs every
10 optimizer steps. Full and ultra run $E=10{,}000$ epochs and validate/log
every 1,000 optimizer steps. As derived in \cref{sec:sampling}, $E$ counts full
passes over the configured sampler and is not generally the optimizer-step
count.

Non-filtering families default to dropping users with an accessible constant
or non-finite lookback in both training and evaluation. The constants family
sets those controls per policy. The reference family exposes independent
overrides for its PatchTST choices. Chronos-2 and persistence set $E=0$ and are
evaluated without optimization.

Evaluation saves elementwise MSE, nMSE, MAE, nMAE, and relative MSE on each
selected split, together with losses grouped by user, equal-user summaries, and
W10 summaries. Tables use \code{test1}, average each metric tensor within a
seed, then report the mean and sample standard deviation across seeds. Each row
uses an explicit $\times10^m$ scale. Every seed directory contains model,
history, metadata, and loss artifacts.

\subsection{Result identities and manifests}
Every family continues as
\code{outputs/family/dataset/L_H/backbone/configs/run_n/seed_n}. The ordered
model-config folders are: \code{policy} for constants;
\code{sampling_mode/batch_size} for sampling; \code{normalization} for
normalizations; \code{loss} for losses; \code{normalization/loss} for the
reference family; \code{normalization} with the linear method as backbone for
linear models; and \code{scope} for central/per-user. No synthetic upper/lower
folder level is required.

Training budget, optimizer, split, strides, and logging/plot controls are
pipeline configs in \code{manifest.json}; scheduler placement is runtime-only.
One seed fixes all stochasticity. \code{EXPERIMENT_MODE} selects path subsets
and is never a phase, family, path component, or computation-signature field.

The current contract uses \code{schema_version: 1} and four states. The default
\code{overwrite_exact} policy skips exact completion, resumes exact
interruption, and allocates a new \code{run_n} for changed pipeline config.
Tables support distinct/latest/selected/average config selection and
selected/latest/distinct/average repeat selection; explicit pipeline filters
must match even with one run. \code{SELECTED_RUNS.txt} stores selection and each
report stores its exact input manifests. Old synchronized outputs lacked the
loss tensors needed for faithful migration and are isolated under
\code{outputs/archive/legacy_pre_schema_v1_2026-08-07}.

The synthetic smoke explicitly enumerates every expected method in all seven
families and requires seeds 1 and 2 for each. Its current run tree is below
\code{outputs/synthetic_smoke/family/synthetic_smoke/24_6/}
\code{numpy_linear_proxy/method/run_n/seed_n}; reports use
\code{outputs/reports/synthetic_smoke}.

\subsection{Main implementation classes}

\begin{table}[htbp]
\centering
\small
\caption{Principal classes and their current responsibilities.}
\label{tab:classes}
\begin{tabularx}{\linewidth}{@{}lX@{}}
\toprule
Class & Responsibility \\
\midrule
\code{TimeSeriesData} & Stores $Z$, dates, stable user IDs, optional covariates, and the target dates owned by a split. \\
\code{IndexSampler} & Builds admissible query dates, applies stride/subsets and pair filtering, and implements the four sampling modes. \\
\code{TimeSeriesDataset} & Materializes $X_{t,u}$, $Y_{t,u}$, covariates, query IDs, user IDs, and window metadata from a sampled pair. \\
\code{TimeTensorModel} & Composes normalization, optional covariate augmentation, a base forecaster, inverse normalization, and optional output rules. \\
\code{LossWrapper} & Applies MSE or MAE after optional lookback-standard-deviation or absolute-lookback-mean scaling. \\
\code{TorchLearner} & Performs one optimizer update per loader batch, records interval-mean training loss, validates at step frequencies, and returns elementwise evaluation losses. \\
\code{SkLinearForecaster} & Deterministically unrolls accessible windows, applies the same reversible normalization contract, and fits/predicts with scikit-learn OLS. \\
\bottomrule
\end{tabularx}
\end{table}

The implemented forward path can therefore be summarized as
\[
  (Z,S,L,H)\xrightarrow{\text{dataset/sampler}}(X,Y)
  \xrightarrow{\text{normalize, forecast, invert}}(\widehat Y,Y)
  \xrightarrow{\text{loss/evaluation}}\text{seed artifacts and tables}.
\]

\end{document}
